% Part: first-order-logic
% Chapter: beyond
% Section: introduction

\documentclass[../../../include/open-logic-section]{subfiles}

\begin{document}

% SEGMENT OLP-0175-S01

\olfileid{fol}{byd}{int}

\olsection{மேலோட்டம்}

ஆர்வமூட்டும் ஒரே தருக்க அமைப்பு முதல்தரத் தருக்கம் அல்ல; முதல்தரத்
தருக்கத்திற்கு பல விரிவாக்கங்களும் மாறுபாடுகளும் உள்ளன. ஒரு தருக்கம்
பொதுவாக ஒரு மொழியின் முறையிய விவரிப்பையும், பெரும்பாலும் (ஆனால்
எப்போதும் அல்ல) ஓர் நிறுவல் அமைப்பையும், பெரும்பாலும் (ஆனால் எப்போதும்
அல்ல) நோக்கமாகக் கொண்ட பொருண்மையியலையும் கொண்டிருக்கும். ஆனால்
இச்சொல்லின் தொழில்நுட்பப் பயன்பாடு இயல்பான கேள்வி ஒன்றை எழுப்புகிறது:
முதல்தரத் தருக்கம் அல்லாத தருக்கங்கள், உள்ளுணர்வு அல்லது மெய்யியல்
பொருளில் பயன்படுத்தப்படும் ``தருக்கம்'' என்ற சொல்லுடன் எவ்வாறு
தொடர்புடையவை? கீழே விவரிக்கப்படும் எல்லா அமைப்புகளும் ஏதோ ஒரு வகை
காரணியத்தை மாதிரிப்பதற்காக வடிவமைக்கப்பட்டவை; அவற்றைத் தருக்கமானவையாக
ஆக்குவது எது என்று கூற முடியுமா?

இதற்கு எளிய விடைகள் இல்லை. ``தருக்கம்'' என்ற சொல் வெவ்வேறு
வழிகளிலும் சூழல்களிலும் பயன்படுத்தப்படுகிறது; ``மெய்மை'' என்ற
கருத்தைப் போலவே, பல மெய்யியல் நிலைப்பாடுகளிலிருந்து அது
பகுப்பாயப்பட்டுள்ளது. எடுத்துக்காட்டாக, எந்தக் கூற்றுகள் அவசியமாக
மெய், அனுபவத்திற்கு முன்னதாகவே மெய், தருக்கமல்லாச் சொற்களின்
பொருள்கோளிலிருந்து சார்பின்றி மெய், அவற்றின் வடிவத்தால் மெய், அல்லது
மொழி மரபால் மெய் என்பதைத் தீர்மானிப்பதே தருக்கக் காரணியத்தின் நோக்கம்
என்று ஒருவர் கொள்ளலாம். இக்கருத்தாக்கங்கள் ஒவ்வொன்றுக்கும் மிகுந்த
தெளிவாக்கம் தேவை. கணிதத்தில் பயன்படுத்தப்படும் தருக்கவகைக்கு மட்டும்
கவனத்தைச் சுருக்கினாலும், அதன் வரம்பு குறித்த ஒப்புமை மிகக் குறைவு.
எடுத்துக்காட்டாக, ஃபிரேகே தொடங்கிய {\em தருக்கவாத} மரபில், ரசலும்
வைட்ஹெட்டும் \textit{பிரின்சிப்பியா மாத்தமாட்டிக்கா} நூலில் தருக்கத்தை
அடிப்படையாகக் கொண்டு கணிதத்தை வளர்க்க முயன்றனர். அவர்களுடைய தருக்க
அமைப்பு கீழே விவரிக்கப்படுவதை ஒத்த உயர்வகைத் தருக்க வடிவமாக இருந்தது.
இறுதியில், பெரும்பாலான அளவுகோல்களின்படி முற்றிலும் தருக்கமானவையாகத்
தோன்றாத அடிகோள்களை (குறிப்பாக முடிவிலி அடிகோளையும் குறைப்புத்தன்மை
அடிகோளையும்) அறிமுகப்படுத்த வேண்டிய நிலை ஏற்பட்டது. இருந்தாலும்
உயர்வரிசைக் காரணியத்தின் சில வடிவங்களைத் தருக்கமானவையாக ஏற்க வேண்டும்
என்று ஒருவர் கொள்ளலாம். மாறாக, ``கூற்றுகளை'' உரையாடலுக்குரிய
சட்டபூர்வப் பொருட்களாக ஏற்காத மெய்ப்பொருளியலைக் கொண்ட குவைன்,
இரண்டாம்வரிசை மற்றும் உயர்வரிசைத் தருக்கங்கள் ஆட்டுத்தோல் போர்த்திய
கணக் கோட்பாடுகளே என்று வாதிடுகிறார். வேறுவிதமாகச் சொன்னால்,
பயனிலைகள் மீதான அளவீட்டைக் கொண்ட அமைப்புகள் முற்றிலும் தருக்கமானவை அல்ல.

இப்போதைக்கு இத்தகைய மெய்யியல் சிக்கல்களைப் பின்னர் கருதுவதற்காக
விட்டுவிட்டு, கீழுள்ள அமைப்புகளைத் தருக்கமானவையோ அல்லாதவையோ ஆன பலவகைக்
காரணியங்களின் முறையியக் கருத்தியலாக்கங்களாக மட்டும் கருதுவது சிறந்தது.

\end{document}
